\chapter{Conception et Dimensionnement de l'Architecture}
\label{chap:conception_dimensionnement}

\section{Segmentation Réseau Complète (VLANs)}
Pour répondre aux exigences de sécurité, d'isolation des flux et de performance de l'institut \textit{DataLab Research}, une segmentation stricte par VLAN (\textit{Virtual Local Area Network}) selon la norme IEEE 802.1Q a été conçue. Cette approche permet de restreindre les domaines de diffusion et d'appliquer des politiques de filtrage strictes entre les différents services et catégories d'utilisateurs.

L'ingénierie réseau de la maquette intègre la totalité des zones requises pour une infrastructure d'entreprise résiliente :
\begin{itemize}
    \item \textbf{VLAN 10 (DMZ - Zone Démilitarisée) :} Zone tampon exposée aux accès externes (futurs serveurs frontaux, proxies).
    \item \textbf{VLAN 20 (Zone Serveurs de Production) :} Zone hautement sécurisée accueillant le cœur des services applicatifs (FTP, DNS, DHCP).
    \item \textbf{VLAN 30 (Réseau Chercheurs) :} Dédié aux stations de travail des 80 scientifiques de l'institut.
    \item \textbf{VLAN 40 (Réseau d'Administration \& Équipe Ingénieurs) :}   Segment isolé  à la maintenance système et pour les postes de développement et d'ingénierie.
    \item \textbf{VLAN 50 (Zone de Supervision) :} Zone réservée à la collecte des métriques.
    \item \textbf{VLAN 60 (Zone de Sauvegarde / Backup) :} Segment d'infrastructure critique dédié à la réplication et à la conservation des données scientifiques.
\end{itemize}

\section{Plan d'Adressage Exhaustif}
Le bloc d'adresses privé global retenu est le réseau \texttt{10.0.0.0/16}. Chaque VLAN se voit attribuer un sous-réseau de classe C (\texttt{/24}), garantissant une lisibilité optimale lors de l'analyse des paquets et un dimensionnement suffisant (jusqu'à 254 hôtes par segment).

Le tableau~\ref{tab:plan_adressage_complet} détaille la cartographie d'adressage IP théorique et validée pour l'ensemble de l'institut.

\begin{table}[htbp]
\centering
\caption{Plan d'adressage complet de l'infrastructure DataLab Research}
\label{tab:plan_adressage_complet}
\begin{tabular}{|l|l|l|l|l|}
\hline
\textbf{VLAN} & \textbf{Nom du Réseau} & \textbf{Plage Réseau} & \textbf{Masque} & \textbf{Passerelle (Gateway)} \\ \hline
VLAN 10       & DMZ Périmétrique       & \texttt{10.0.10.0/24}  & \texttt{255.255.255.0} & \texttt{10.0.10.254}       \\ \hline
VLAN 20       & Serveurs Production    & \texttt{10.0.20.0/24}  & \texttt{255.255.255.0} & \texttt{10.0.20.254}       \\ \hline
VLAN 30       & Zone Chercheurs        & \texttt{10.0.30.0/24}  & \texttt{255.255.255.0} & \texttt{10.0.30.254}       \\ \hline
VLAN 40       & Admin \& Équipe Ingénieurs      & \texttt{10.0.40.0/24}  & \texttt{255.255.255.0} & \texttt{10.0.40.254}       \\ \hline
VLAN 50       & Supervision   & \texttt{10.0.50.0/24}  & \texttt{255.255.255.0} & \texttt{10.0.50.254}       \\ \hline
VLAN 60       & Stockage \& Backup     & \texttt{10.0.60.0/24}  & \texttt{255.255.255.0} & \texttt{10.0.60.254}       \\ \hline
\end{tabular}
\end{table}


La figure \ref{fig:vision} illustre la vision globale de l'architecture réseau et logique de la maquette.
\begin{figure}[htbp]
	\centering
	\includegraphics[width=0.85\textwidth]{figures/vision_globale.jpeg}
	\caption{vision globale de l'architecture réseau et logique de la maquette}
	\label{fig:vision}
\end{figure}


% \section{Topologie Générale de la Maquette}
% La centralisation du routage inter-VLAN est opérée au niveau du cœur de réseau selon un modèle \textit{Router-on-a-Stick}. Le routeur principal de la maquette s'interface avec le commutateur au moyen d'un lien d'agrégation (\textit{Trunk}), où chaque sous-interface virtuelle du routeur gère l'encapsulation d'un VLAN spécifique.

% La figure~\ref{fig:archi_globale} illustre la topologie globale et l'interconnexion de nos serveurs virtuels au sein du simulateur.

% \begin{figure}[htbp]
%     \centering
%     \includegraphics[width=0.85\textwidth]{figures/archi_globale.jpeg}
%     \caption{Topologie réseau et architecture logique sous GNS3}
%     \label{fig:archi_globale}
% \end{figure}

